\section{Check access histories as well as results}
The physical ring is 32 bytes long, but only the first \(C\) bytes belong to a
selected descriptor. A word load crossing from index \(C-1\) into mapped padding
is outside its logical allocation. The source pointer similarly authorizes only
its requested payload length for push. A zero-length push reads no source bytes,
even though it still stores a length prefix in the ring.

The checks inspect native addresses and full access widths. They do not require
each input operation to be a byte load. A safely aligned wider operation can be
valid if its entire width respects the relevant interval and every written byte
belongs to that successful operation's permitted destination. Reading or writing
mapped padding is not justified by a correct final queue.

Failure checks also retain store histories. An implementation that decrements
used before checking destination capacity and later restores it has the same final
state as the input, but violates the no-write failure rule. Pop must not clear old
ring bytes after copying them, because this contract explicitly preserves ring
contents. On success, bytes beyond the returned payload in the destination remain
unchanged and receive no stores.

Each independent call test changes descriptor selection, ring selection, source
and destination offsets, full saved-register
canaries. Queue states cover empty and full cases, exact fits, zero-length records,
and wrapped data. No test requires the example's scratch registers, payload store
order, or instruction count. The cumulative lesson checks separately inspect
particular intermediate registers to verify what the text illustrates.
